\DocumentMetadata{
  pdfversion=2.0,
  pdfstandard=ua-2,
  lang=en-US,
  tagging=on,
  tagging-setup={math/setup=mathml-SE,tabsorder=structure}
}
\documentclass[11pt,letterpaper]{article}
\usepackage[margin=0.68in,includefoot]{geometry}
\usepackage{newcomputermodern}
\usepackage{amsmath,amscd,mathtools}
\usepackage{tikz,adjustbox}
\providecommand{\square}{\mdlgwhtsquare}
\usepackage[hidelinks]{hyperref}
\hypersetup{
  pdftitle={Algebra PhD Exam, September 1991},
  pdfauthor={University of Florida Department of Mathematics},
  pdfsubject={Graduate mathematics examination},
  pdfdisplaydoctitle=true,
  bookmarksopen=true
}
\setcounter{secnumdepth}{0}
\setlength{\parindent}{0pt}
\setlength{\parskip}{0.28em}
\setlength{\emergencystretch}{3em}
\setlength{\leftmargini}{2.15em}
\setlength{\leftmarginii}{2.65em}
\setlength{\leftmarginiii}{3em}
\pagestyle{empty}
\raggedbottom
\allowdisplaybreaks
\definecolor{ProblemConcern}{HTML}{7A1F1F}
\newenvironment{problemconcern}{%
  \par\tagstructbegin{tag=Aside}\begingroup\color{ProblemConcern}
}{%
  \par\endgroup\tagstructend\nopagebreak[4]\vspace{0.12em}
}
\NewTaggingSocketPlug{block/list/label}{examproblem}{%
  \tagstructbegin{tag=Lbl}\tagstructend%
  \tagstructbegin{tag=\LBody}%
  \tagstructbegin{tag=\ProblemHeadingTag,title-o={\ProblemHeadingText}}%
  \tagmcbegin{tag=\ProblemHeadingTag,actualtext={\ProblemHeadingText}}%
  #2%
  \tagmcend%
  \tagstructend%
}
\newenvironment{examproblems}[2]{%
  \def\ProblemHeadingTag{#2}%
  \AssignTaggingSocketPlug{block/list/label}{examproblem}%
  \begin{enumerate}%
  \setcounter{enumi}{#1}%
  \renewcommand{\labelenumi}{\textbf{\theenumi.}}%
  \setlength{\topsep}{0.35em}%
  \setlength{\partopsep}{0pt}%
  \setlength{\itemsep}{0.48em}%
  \setlength{\parsep}{0.18em}%
}{\end{enumerate}}
\newenvironment{examsubparts}{%
  \AssignTaggingSocketPlug{block/list/label}{default}%
  \begin{enumerate}%
  \setlength{\topsep}{0.18em}%
  \setlength{\partopsep}{0pt}%
  \setlength{\itemsep}{0.16em}%
  \setlength{\parsep}{0.1em}%
}{\end{enumerate}}
\newenvironment{examinstructions}{%
  \par\begingroup\itshape
}{%
  \par\endgroup\vspace{0.18em}
}
\makeatletter
\renewcommand\subsection{\@startsection{subsection}{2}{0pt}%
  {-1.25ex plus -.25ex minus -.1ex}%
  {0.55ex plus .1ex}%
  {\normalfont\large\bfseries}}
\renewcommand{\@maketitle}{%
  \newpage\null\vskip -1em%
  {\footnotesize\bfseries
    \href{https://www.ufl.edu/}{University of Florida}, \href{https://math.ufl.edu/}{Department of Mathematics}\par}%
  \vskip -0.5em%
  \tagstructbegin{tag=H1,title={Algebra PhD Exam, September 1991}}%
  \tagpdfparaOff%
  \tagmcbegin{tag=H1}%
    {\LARGE\bfseries\href{https://gma.math.ufl.edu/exams/algebra-phd/}{Algebra PhD Exam}, September 1991\par}%
  \tagmcend%
  \tagpdfparaOn%
  \tagstructend%
  \vskip 0.25em%
  \tagmcbegin{artifact}%
    \hrule height 1pt%
  \tagmcend%
  \vskip 1em%
}
\makeatother
\title{Algebra PhD Exam, September 1991}
\author{}
\date{}
\begin{document}
\maketitle
\thispagestyle{empty}
\begin{examinstructions}
Please read through all the questions before you begin. Your seven best solutions will count. You may quote standard results (within reason) as long as you make it clear that you are doing so and state them clearly.
\end{examinstructions}
\begin{examproblems}{0}{H2}
\def\ProblemHeadingText{Problem 1.}
\item
Prove that, if~\(F\) is a field, the group of upper triangular~\(n\) by~\(n\) matrices (\(n>1\)) with entries in~\(F\) and 1 on every diagonal entry is nilpotent of class \(n-1\).
\def\ProblemHeadingText{Problem 2.}
\item
Calculate in detail the Galois group of the polynomial \(x^3+2\) over the field of rational numbers.
\def\ProblemHeadingText{Problem 3.}
\item
\begin{problemconcern}
\textbf{Warning: Suspected error.} The source clearly appears to say “upper triangular 2 by 2 matrices with integer coefficients,” but as written these matrices do not specify the standard free nilpotent class 2 group on two generators. A possible intended formulation involving upper unitriangular 3 by 3 integer matrices is not uniquely determined from the source, so the source reading has been retained.
\end{problemconcern}
Prove that the group~\(G\) of upper triangular 2 by 2 matrices with integer coefficients is a free nilpotent class 2 group on two generators.
\def\ProblemHeadingText{Problem 4.}
\item
Suppose~\(R\) is a finite commutative ring with identity in which \(r^3=r\) for all \(r\in R\). Prove that~\(R\) is a finite direct sum of fields, each of which is isomorphic to \(\mathbb Z_2\) or \(\mathbb Z_3\).
\def\ProblemHeadingText{Problem 5.}
\item
State and prove Hilbert’s Basis Theorem. Let \(R=\mathbb Z[x_1,x_2,\ldots]\) be the ring of polynomials in countably many independent variables with integer coefficients. Is~\(R\) Noetherian? Justify your answer.
\def\ProblemHeadingText{Problem 6.}
\item
Let~\(G\) be a group and~\(H\) be a normal subgroup of~\(G\), and assume that the center of~\(H\) is the identity. Show that the following are equivalent. Give an example to show that if~\(H\) had a non-trivial center the two statements might not be equivalent.
\begin{examsubparts}
\item[\textnormal{a)}]
There is a subgroup~\(J\) of~\(G\) such that~\(G\) is the (internal) direct product of~\(H\) and~\(J\).
\item[\textnormal{b)}]
For every \(g\in G\), \(g\) induces by conjugation an inner automorphism of~\(H\).
\end{examsubparts}
\def\ProblemHeadingText{Problem 7.}
\item
Let~\(R\) be a commutative ring with 1 and~\(S\) a subset of non-zero divisors of~\(R\) such that \(0\notin S\) and~\(S\) is a semigroup. Form the ring of quotients \(S^{-1}R\). Prove that the prime ideals of \(S^{-1}R\) are in one-to-one correspondence with the prime ideals~\(P\) of~\(R\) such that \(P\cap S=\varnothing\).
\def\ProblemHeadingText{Problem 8.}
\item
If~\(R\) is a commutative ring with identity prove that
\[
\bigcap\{P\mid P\text{ is a prime ideal of }R\}=\{x\in R\mid x^n=0\text{ for some }n>0\}.
\]
 (use Zorn’s Lemma).
\def\ProblemHeadingText{Problem 9.}
\item
Prove that if~\(F\) is a free abelian group then each non-trivial subgroup~\(H\) of~\(F\) is free abelian. (Hint: \(F=\) direct sum of \(\mathbb Z\)’s. Well order the copies. Define \(F(i)\) to be the sum of all the copies of \(\mathbb Z\) assigned some \(j<i\). Then \(F(j+1)/F(j)\cong\mathbb Z\) and the union of all the \(F(i)\) is~\(F\). Now consider \(H(i)=H\cap F(i)\). Prove that \(H(i+1)/H(i)\) is isomorphic either to \(\mathbb Z\) or to \(\{0\}\) and that the union of all the \(H(i)\) is~\(H\).)
\def\ProblemHeadingText{Problem 10.}
\item
Prove that if~\(k\) is a field then the power series ring \(k[[x]]\) in one variable is a discrete valuation ring.
\def\ProblemHeadingText{Problem 11.}
\item
Let~\(R\) be a ring with identity and let~\(A\) be a right-\(R\)-module and~\(B\) a left-\(R\)-module. Define \(A\otimes_R B\). Let
\[
0\longrightarrow N\longrightarrow M\longrightarrow L\longrightarrow 0
\]
 be an exact sequence of left-\(R\)-modules. Show that
\[
A\otimes_R N\longrightarrow A\otimes_R M\longrightarrow A\otimes_R L\longrightarrow 0
\]
 is exact. Give an example to show that the first map need not be injective.
\end{examproblems}
\end{document}
